% Generated draft: requires mathematical and source review.
\documentclass[11pt]{article}
\usepackage[margin=25mm]{geometry}
\usepackage{fontspec}
\setmainfont{Latin Modern Roman}
\usepackage{amsmath,amssymb,mathtools,unicode-math}
\setmathfont{Latin Modern Math}
\usepackage{enumitem,needspace,xurl,hyperref}
\setlength{\parindent}{0pt}
\setlength{\parskip}{5pt}
\setlength{\emergencystretch}{3em}
\newcommand{\N}{\mathbb{N}}
\newcommand{\Z}{\mathbb{Z}}
\newcommand{\Q}{\mathbb{Q}}
\newcommand{\R}{\mathbb{R}}
\newcommand{\C}{\mathbb{C}}
\newcommand{\F}{\mathbb{F}}
\newcommand{\A}{\mathbb{A}}
\newcommand{\PP}{\mathbb{P}}
\newcommand{\E}{\mathbb{E}}
\newcommand{\eps}{\varepsilon}
\newcommand{\dd}{\,\mathrm{d}}
\begin{document}
\section*{UnsolvedMath ID: 30002054}
\textbf{Source problem code:} catalogue entry
\subsection{Definitions and mathematical statement}
Fix an integer \(d\geq 4\). For a finite simplicial triangulation \(\Delta\) of a compact \((d-1)\)-dimensional manifold \(M\), let \(f_j(\Delta)\) denote its number of \(j\)-dimensional faces, with \(f_{-1}(\Delta)=1\). Define its \(h\)-vector by
\[
\sum_{i=0}^{d}h_i(\Delta)t^{d-i}
=\sum_{i=0}^{d}f_{i-1}(\Delta)(t-1)^{d-i}.
\]
Let \(f_0^{\mathrm{int}}(\Delta)\) be the number of vertices whose images lie in \(M\setminus\partial M\). Write \(|\Delta|\cong M\) for a triangulation whose geometric realization is homeomorphic to \(M\).

For a manifold with boundary, the supplied record gives the definition
\[
\Gamma(M)
=\min_{|\Delta|\cong M}
\bigl(h_2(\Delta)-h_{d-1}(\Delta)-d\,h_d(\Delta)\bigr)
=\min_{|\Delta|\cong M}
\bigl(h_2(\Delta)-f_0^{\mathrm{int}}(\Delta)\bigr).
\]
The equality of these two minimum formulas is part of the supplied formulation. The record does not explain why it holds for the entire stated class of compact manifolds with boundary, or identify any additional hypotheses needed for it. This draft therefore retains the asserted equality as a source qualification requiring review, rather than establishing that the two expressions define the same invariant. The record also does not specify a category of triangulations or hypotheses ensuring that the triangulation class is nonempty and that the minima exist as finite values.

For a closed manifold, the supplied extension is
\[
\Gamma(M)=\Gamma(M\setminus B^{d-1}).
\]
Here the notation for the removed ball is unresolved. To obtain a compact manifold with boundary, the customary puncturing construction removes the interior of an embedded closed \((d-1)\)-ball. The record does not explicitly specify this convention, the embedding category, or independence of the choice of ball.

\textbf{Supplied additivity question.} For compact manifolds \(M_1,M_2\) of the same dimension \(d-1\), with \(d\geq4\), does
\[
\Gamma(M_1\mathbin{\#}M_2)
=\Gamma(M_1)+\Gamma(M_2)
\]
hold for every admissible pair? The intended class of admissible pairs and the meaning of \(\#\) require clarification: the record does not specify whether the target concerns closed manifolds, manifolds with boundary, or both, nor whether a boundary case uses interior connected sum or boundary connected sum. Connectedness, orientability, and gluing conventions are likewise unstated. Thus the universal target is preserved, but its precise domain and the well-definedness of \(\Gamma\) remain subject to the qualifications above.
\subsection{Short English statement}
The supplied problem asks whether a minimum triangulation quantity \(\Gamma\) is additive under connected sum in dimensions at least three. It gives two asserted formulas for \(\Gamma\) and extends it to closed manifolds by puncturing. The equivalence of the formulas, admissible triangulations, puncturing convention, and scope of connected sum require clarification. The supplied open status is dated curation, not an independently verified conclusion.
\subsection{Sources}
\begin{itemize}
\item Local UnsolvedMath registry: supplied complete source record, permanent ID \(30002054\), exact source problem code \texttt{catalogue entry}, title \emph{Additivity of Triangulation Complexity under Connected Sum}. Registry paths identified by the project are \texttt{data/unsolvedmath/problems.json} and \texttt{catalog/problems.json}; this draft uses the supplied record without independently inspecting those files.
\item Supplied primary citation: \emph{Triangulations (2013)}, Oberwolfach Reports, pp.~1405--1486, DOI \texttt{10.4171/owr/2012/24}. \url{https://doi.org/10.4171/owr/2012/24}. The record attributes the problem to the report's workshop open-problem discussions, gives proposed year 2012, and names no individual proposer. The relationship between the supplied 2013 citation year and the 2012 identifiers requires bibliographic review.
\item Supplied literature source: Oberwolfach report, \emph{Triangulations}. \url{https://doi.org/10.4171/owr/2012/24}.
\item Supplied literature source: De Loera--Rambau--Santos, \emph{Triangulations}. \url{https://link.springer.com/book/10.1007/978-3-642-12971-1}.
\item Supplied literature assessment, checked on \texttt{2026-08-21}: the record labels the problem open and reports that no proof or counterexample for the stated invariant was located. This is supplied dated curation derived from the OWR material, not an independently verified literature assessment; the record itself requests independent verification of time-sensitive status and cited claims.
\item The supplied statement status is \texttt{corrected\_verified}, with a claim that the clean statement was checked against the official report. This is preserved as source metadata; this draft does not claim to have visited or verified the report.
\end{itemize}
\end{document}
